% !TeX program = xelatex
\documentclass[UTF8,zihao=-4]{ctexart}
% XeLaTeX 下显式选择可嵌入的 CJK 字体。某些 PDF 预览器无法解析
% ctex 默认 Fandol 字体生成的 Adobe-GB1 CMap，导致中文（包括图中的节点文字）空白。
% 若运行环境没有 Noto CJK，则保留 ctex 的默认字体回退。
\IfFontExistsTF{Noto Serif CJK SC}{%
  \setCJKmainfont[AutoFakeSlant=true]{Noto Serif CJK SC}%
}{}
\IfFontExistsTF{Noto Sans CJK SC}{%
  \setCJKsansfont{Noto Sans CJK SC}%
}{}
\IfFontExistsTF{Noto Sans Mono CJK SC}{%
  \setCJKmonofont{Noto Sans Mono CJK SC}%
}{}
\usepackage[a4paper,margin=2.15cm]{geometry}
\usepackage{amsmath,amssymb,mathtools,bm}
\usepackage{booktabs,longtable,tabularx,array}
\usepackage{xcolor}
\usepackage{listings}
\usepackage{float,placeins}
\usepackage{algorithm}
\usepackage{algpseudocode}
\usepackage{tikz}
\usetikzlibrary{arrows.meta,positioning,shapes.geometric,fit,calc}
\usepackage{hyperref}
\hypersetup{hidelinks}

\definecolor{iceblue}{HTML}{EAF3F8}
\definecolor{waterblue}{HTML}{DCEFF8}
\definecolor{heatred}{HTML}{FCE8E6}
\definecolor{codebg}{HTML}{F6F8FA}
\definecolor{codeframe}{HTML}{D0D7DE}
\definecolor{codekeyword}{HTML}{CF222E}
\definecolor{codecomment}{HTML}{57606A}
\lstdefinestyle{iceflow}{
  language=Python,
  basicstyle=\ttfamily\small,
  keywordstyle=\color{codekeyword}\bfseries,
  commentstyle=\color{codecomment}\itshape,
  backgroundcolor=\color{codebg},
  frame=single,
  rulecolor=\color{codeframe},
  framerule=0.6pt,
  framesep=8pt,
  columns=fullflexible,
  keepspaces=true,
  showstringspaces=false,
  breaklines=true,
  postbreak=\mbox{\textcolor{gray}{$\hookrightarrow$}\space}
}

\floatname{algorithm}{算法}
\algrenewcommand\algorithmicrequire{\textbf{输入：}}
\algrenewcommand\algorithmicensure{\textbf{输出：}}
\algrenewcommand\algorithmicif{\textbf{若}}
\algrenewcommand\algorithmicthen{\textbf{则}}
\algrenewcommand\algorithmicelse{\textbf{否则}}
\algrenewcommand\algorithmicfor{\textbf{对}}
\algrenewcommand\algorithmicdo{\textbf{执行}}
\algrenewcommand\algorithmicend{\textbf{结束}}

\newcolumntype{Y}{>{\raggedright\arraybackslash}X}
\newcommand{\vect}[1]{\boldsymbol{#1}}
\newcommand{\dd}{\mathop{}\!\mathrm d}
\newcommand{\clip}{\operatorname{clip}}
\newcommand{\sgn}{\operatorname{sgn}}
\newcommand{\code}[1]{\texttt{#1}}
\newcommand{\cellarea}{A_c}
\newcommand{\stepname}[1]{\paragraph{对应代码：\code{#1}.}}

\tikzset{
  flow/.style={draw,rounded corners=2pt,align=center,minimum height=8mm,
    text width=10.8cm,fill=iceblue,font=\small},
  slow/.style={flow,fill=heatred},
  decision/.style={diamond,draw,aspect=3.0,align=center,fill=white,
    inner sep=1.5pt,font=\small},
  branch/.style={draw,rounded corners=2pt,align=center,minimum height=11mm,
    text width=3.8cm,fill=waterblue,font=\small},
  arrow/.style={-{Latex[length=2.2mm]},thick}
}

\numberwithin{equation}{section}
\setcounter{tocdepth}{3}
\setcounter{secnumdepth}{3}
\setlength{\emergencystretch}{3em}
\setlength{\parskip}{0.25em}

\title{IceFlow2D：下落冰块--热水--融化二维耦合方法\\
\large 从物理概念到 \code{step()} 每一行的离散实现}
\author{}
\date{对应当前 \code{iceflow2d/simulator.py} 与 \code{iceflow2d/thermal.py}}

\begin{document}
\maketitle

\begin{abstract}
本文说明 IceFlow2D 中“冰块从空气落入热水、受流体作用平移和转动、受容器边界防穿墙约束并逐渐融化”
这一专用二维求解器。正文面向第一次接触格子 Boltzmann 方法、相场法、任意拉格朗日--欧拉
重映射和刚体约束的新读者：先解释每个数组究竟保存什么，再给出连续模型意图、格子单位转换、
D2Q9 离散、初始化方法，最后严格按照 \code{IceFlow2D.step()} 的调用顺序逐项展开。
特别需要强调的是，程序并不是用一条统一的“焓方程”同时搬运冰与水：固冰质量和显热保存在
随体材料网格，液水面积和显热保存在固定世界网格；输出中的体积焓只是由这些守恒广延量合成的
诊断场。本文所有二维质量、体积、能量和惯量均按单位出平面厚度解释。
\end{abstract}

\tableofcontents
\clearpage

\section{阅读路线、范围与最小直觉}

\subsection{本程序在解什么问题}

计算域是一只二维容器。下部是热水，上部是空气，中间存在可变形的弥散自由液面；一个初始为
刚性方块的冰体从空气区下落，穿过自由液面后继续受到水动力和重力；到达容器边界时只施加
防止整体穿墙的位置与法向速度限制。冰在材料
坐标中导热并融化；剩余固冰仍作为一个刚体运动，融水则转入世界坐标中的液水，并最终进入空气--
水相场和流体动量。

求解器当前只保留这一条专用路径：
\begin{itemize}
  \item 空气--水：压力--动量 D2Q9 LBM 与守恒 Allen--Cahn 相场；
  \item 冰--流体：锐界 cut-link 反弹与逐链动量交换；
  \item 冰体：动态质量、质心和惯量的单刚体积分；
  \item 热学：材料坐标固冰显热、世界坐标液水显热、保守 ALE 重映射；
  \item 相变：显热先把冰升至熔点，再支付潜热生成融水；
  \item 守恒闭合：相场水量、热水 aperture、融水质量、线动量和角动量分别闭合。
\end{itemize}
程序不包含断裂、多刚体、蒸发、沸腾、蒸汽输运、气泡、再冻结，也不把黏性耗散或墙面速度截断
造成的机械能变化转成热。

\subsection{六个初学者容易混淆的词}

\begin{description}
  \item[格点与格子单元] LBM 的未知量放在规则网格点（代码中也以单元中心位置
  $(i+\tfrac12,j+\tfrac12)$ 计算几何）。本文把索引位置统称为世界单元 $P=(i,j)$。
  \item[population（分布函数）] 每个世界单元不是只存一个速度，而是沿九个离散方向各存
  一个数。碰撞改变本地九个数，流送把它们沿各方向送到邻格。
  \item[相场 $\phi$] $\phi\approx1$ 表示水，$\phi\approx0$ 表示空气，二者之间的连续过渡带
  是自由液面。它不是“冰的液相率”。
  \item[锐界冰掩膜] 冰的连续材料覆盖率经过阈值化后得到 $\chi^s\in\{0,1\}$，决定一个格点
  是否参加 LBM。连续覆盖率仍用于热接触，不能与锐掩膜混为一谈。
  \item[aperture] 一个世界单元能够容纳液水的有效开口或容量。热模块中的液水面积必须放在
  LBM 认为合法的水侧开口中。
  \item[ALE 与 refill] 刚体移动会扫过或释放世界单元。ALE 重映射守恒搬运热水广延量；
  refill 则为新开放的 LBM 节点重建 $f_q,h_q$。两者服务于不同未知量。
\end{description}

\subsection{连续模型与离散实现的关系}

后文会给出连续方程，帮助读者理解“算法想逼近什么”；但数值结果真正由离散公式决定。
因此，本文把“目标方程”和“代码精确离散”分开陈述，不由源码无法证明的内容推断全局收敛阶、
无条件稳定性或严格解析守恒。

\section{几何、坐标、单位和默认工况}

\subsection{两个坐标系与两个刚体原点}

世界坐标 $\vect x$ 固定在容器上；材料坐标 $\vect\xi$ 固定在初始冰体上。材料参考原点记为
$\vect X_0$，当前剩余固冰的材料质心记为 $\vect\xi_c$，世界质心记为 $\vect X_c$。刚体映射为
\begin{equation}
  \vect x(\vect\xi,t)=\vect X_0(t)+\vect R(\theta)\vect\xi,
  \qquad
  \vect X_c=\vect X_0+\vect R(\theta)\vect\xi_c,
  \label{eq:rigid-map}
\end{equation}
其中
\begin{equation}
  \vect R(\theta)=
  \begin{bmatrix}\cos\theta&-\sin\theta\\ \sin\theta&\cos\theta\end{bmatrix}.
\end{equation}
代码字段 \code{body\_reference\_origin} 对应 $\vect X_0$，
\code{body\_center} 对应 $\vect X_c$。非对称融化会改变 $\vect\xi_c$，所以二者一般不重合。
任意冰面点的格子速度为
\begin{equation}
  \vect u_b(\vect x)=\vect V_c+\omega\,\vect e_z\times(\vect x-\vect X_c),
  \qquad
  \vect e_z\times(r_x,r_y)=(-r_y,r_x).
  \label{eq:body-point-velocity}
\end{equation}

\subsection{物理量与格子量}

LBM 在格子单位中取 $\Delta x_{\rm LU}=\Delta t_{\rm LU}=1$。程序固定令物理参考速度
$U_{\rm ref}$ 对应 $u_{\rm ref,LU}=0.1$，于是
\begin{equation}
  \Delta t=\frac{0.1\,\Delta x}{U_{\rm ref}},
  \qquad
  \vect u_{\rm LU}=\vect u_{\rm phys}\frac{\Delta t}{\Delta x}.
  \label{eq:lattice-time-map}
\end{equation}
物理运动黏度、重力加速度和表面张力映射为
\begin{align}
  \nu_{\rm LU}&=\nu_{\rm phys}\frac{\Delta t}{\Delta x^2},\\
  \vect g_{\rm LU}&=\vect g_{\rm phys}\frac{\Delta t^2}{\Delta x},\\
  \sigma_{\rm LU}&=\sigma_{\rm phys}
    \frac{(0.1)^2}{\rho_wU_{\rm ref}^2\Delta x}.
  \label{eq:physical-lattice-scales}
\end{align}
LBM 密度以水密度归一化：$\rho_{w,\rm LU}=1$、
$\rho_{a,\rm LU}=\rho_a/\rho_w$。刚体材料单元的物理二维质量除以
$\rho_w\Delta x^2$ 后进入格子动力学。

\subsection{命令行示例的当前默认值}

表~\ref{tab:cli-defaults} 对应
\path{iceflow2d/examples/coupled_falling_ice_melting_2d.py} 的无参数运行。
这与直接构造 \code{IceFlowConfig()} 的类默认值不同：后者仍为 $200\times400$、
$\Delta x=1.25\times10^{-4}\,\mathrm m$。本文讨论用户实际运行入口，因此以下数值采用 CLI 默认值。

\begin{table}[H]
  \centering
  \caption{当前命令行示例的默认下落融化工况}
  \label{tab:cli-defaults}
  \begin{tabularx}{\textwidth}{@{}lY@{}}
    \toprule
    项目 & 默认值与含义 \\
    \midrule
    区域与网格 & $0.025\,\mathrm m\times0.050\,\mathrm m$，$100\times200$，正方形单元 \\
    网格与时间步 & $\Delta x=2.5\times10^{-4}\,\mathrm m$，
      $\Delta t=6.25\times10^{-6}\,\mathrm s$ \\
    水池与冰块 & 水深 $0.030\,\mathrm m=120$ 格；方冰边长 $0.008\,\mathrm m=32$ 格 \\
    初始位置 & 冰块旋转 $5^\circ$；最低角点高于水面 $0.0015\,\mathrm m=6$ 格 \\
    初始温度 & 水/冰/空气 $=90/0/20\,^{\circ}\mathrm C$；熔点 $T_m=0\,^{\circ}\mathrm C$ \\
    密度 & $\rho_w=1000$、$\rho_i=917$、$\rho_a=1.25\,\mathrm{kg\,m^{-3}}$ \\
    黏度与表面张力 & $\nu_w=10^{-6}$、$\nu_a=1.5\times10^{-5}\,\mathrm{m^2s^{-1}}$；
      $\sigma=0.072\,\mathrm{N\,m^{-1}}$ \\
    重力与速度映射 & $\vect g=(0,-9.8)\,\mathrm{m\,s^{-2}}$；
      $4\,\mathrm{m\,s^{-1}}\leftrightarrow0.1$ 格/步 \\
    相场 & 界面宽度 $W=5$ 格，迁移率 $M=0.1$，预热流送 500 步 \\
    热学 & 慢热间隔 $N_{\rm slow}=8$；$Fo_{\max}=0.15$，$Co_{\max}=0.50$ \\
    热边界 & 四面容器墙绝热；水--气与冰--气绝热 \\
    输出 & 总时长 $3\,\mathrm s$，每 $0.01\,\mathrm s$ 采样 \\
    \bottomrule
  \end{tabularx}
\end{table}

因此一个慢热区间为
\begin{equation}
  \Delta t_{\rm slow}=8\Delta t=5.0\times10^{-5}\,\mathrm s,
\end{equation}
每个输出间隔含 1600 个 LBM 步，终点为 480000 步；若计入 $t=0$，共 301 个采样。

\section{未知量、时间层和符号表}

\subsection{世界流体场}

每个活动世界单元 $P$ 保存九个压力--动量分布 $f_{Pq}$、九个相场分布 $h_{Pq}$，以及
$\phi_P,\vect u_P,p_P,\vect F_P$。定义墙掩膜 $\chi_P^w$ 与锐冰掩膜 $\chi_P^s$，活动集合为
\begin{equation}
  \mathcal A^n=\{P\mid \chi_P^w=0,\ \chi_P^{s,n}=0\}.
  \label{eq:active-set}
\end{equation}
这里 \code{fluid\_force} 对应的 $\vect F$ 是格子加速度，即力密度除以局部混合密度后的量。
流送恢复的存储速度与实际用于平衡态、压力和热对流的半力速度分别为
\begin{equation}
  \vect u=\frac{\sum_q\vect c_qf_q}{\rho},
  \qquad
  \vect u^\star=\vect u+\frac12\vect F.
  \label{eq:half-force-velocity}
\end{equation}
这一差别非常重要：输出和物理解释应使用 $\vect u^\star$。

\subsection{材料冰和世界水的热广延量}

材料单元用希腊索引 $\alpha$，保存初始固冰质量 $m_\alpha^0$、当前固冰质量
$m_\alpha^s$ 和相对熔点的固冰显热 $E_\alpha^s$：
\begin{equation}
  f_\alpha^s=\frac{m_\alpha^s}{m_\alpha^0},
  \qquad
  T_{i,\alpha}=T_m+\frac{E_\alpha^s}{m_\alpha^sc_{p,i}}\le T_m.
  \label{eq:body-thermal-state}
\end{equation}
世界单元保存液水面积 $V_P^w$ 与相对熔点的液水显热 $S_P^w$：
\begin{equation}
  S_P^w=\rho_wc_{p,w}V_P^w(T_{w,P}-T_m),
  \qquad 0\le V_P^w\le \cellarea=\Delta x^2.
  \label{eq:water-extensive-state}
\end{equation}
由于是二维单位厚度模型，$m_\alpha^s$ 的单位为 $\mathrm{kg\,m^{-1}}$，
$E_\alpha^s,S_P^w$ 的单位为 $\mathrm{J\,m^{-1}}$，$V_P^w$ 的单位为 $\mathrm m^2$。

输出场 \code{enthalpy\_j\_m3} 不是推进未知量，而是诊断合成：
\begin{equation}
  H_P=
  \begin{cases}
  (\rho_wLV_P^w+S_P^w)/\cellarea,&P\text{ 显示为水},\\
  [E_\alpha^s+(m_\alpha^0-m_\alpha^s)L]/\cellarea,&P\text{ 显示为材料冰}.
  \end{cases}
  \label{eq:diagnostic-enthalpy}
\end{equation}

\subsection{常用符号速查}

\begin{longtable}{@{}p{0.19\textwidth}p{0.27\textwidth}p{0.46\textwidth}@{}}
\caption{数学符号、代码含义及存储位置}\label{tab:symbols}\\
\toprule
符号 & 主要代码字段 & 含义 \\
\midrule
\endfirsthead
\toprule
符号 & 主要代码字段 & 含义 \\
\midrule
\endhead
$P=(i,j)$，$q$ & 数组索引 & 世界单元及 D2Q9 方向编号 \\
$\phi$ & \code{phi} & 空气--水相场；水约为 1，空气约为 0 \\
$f_q,h_q$ & \code{f}, \code{h} & 压力--动量分布与相场分布 \\
$\vect u,\vect u^\star$ & \code{u}, \code{u + force/2} & 存储速度与半力物理速度 \\
$p^H,p^d,p$ & hydrostatic/reference, \code{p} & 静水参考、动态压力及总压力，$p=p^H+p^d$ \\
$C_P$ & \code{world\_body\_indicator} & 当前位姿下连续冰覆盖率 \\
$\chi_P^s$ & \code{solid} & 由 $C_P$ 阈值化得到的锐冰掩膜 \\
$d_P$ & \code{sdf} & 冰界有符号距离；共享世界 SDF 的物理单位为米 \\
$\vect X_0,\vect X_c$ & reference origin, body center & 材料参考原点与当前质量中心 \\
$\vect V_c,\omega$ & body velocity/angular velocity & 质心格子速度与每步转角 \\
$m_b,I_b$ & body mass/inertia & 由剩余材料质量归约的格子质量与惯量 \\
$V_P^w,S_P^w$ & thermal water volume/sensible & 世界水面积与水显热广延量 \\
$m_\alpha^s,E_\alpha^s$ & body solid mass/sensible & 材料固冰质量与固冰显热广延量 \\
$N_{\rm slow}$ & update interval & 两次慢热更新之间的 LBM 步数 \\
\bottomrule
\end{longtable}

\section{D2Q9 与连续模型意图}

\subsection{九个离散速度}

D2Q9 的离散速度和权重为
\begin{align}
  \vect c_0&=(0,0),\\
  \vect c_{1\ldots4}&=(1,0),(0,1),(-1,0),(0,-1),\\
  \vect c_{5\ldots8}&=(1,1),(-1,1),(-1,-1),(1,-1),\\
  w_0&=\frac49,\qquad w_{1\ldots4}=\frac19,
  \qquad w_{5\ldots8}=\frac1{36},\qquad c_s^2=\frac13.
  \label{eq:d2q9}
\end{align}
权重满足 $\sum_qw_q=1$、$\sum_qw_q\vect c_q=\vect0$ 和
$\sum_qw_q\vect c_q\vect c_q=c_s^2\vect I$，这些恒等式使宏观矩可以从九个 population 恢复。

\subsection{空气--水目标方程}

相场混合密度与运动黏度分别取
\begin{align}
  \bar\phi&=\clip(\phi,0,1),\\
  \rho(\phi)&=\rho_a+(\rho_w-\rho_a)\bar\phi,\\
  \nu_0(\phi)&=\nu_a+(\nu_w-\nu_a)\bar\phi.
  \label{eq:mixture-properties}
\end{align}
相场部分意图逼近守恒 Allen--Cahn 方程
\begin{equation}
  \partial_t\phi+\nabla\cdot(\phi\vect u^\star)
  =\nabla\cdot\left\{M\left[\nabla\phi-\Theta(\phi)\vect n\right]\right\},
  \quad
  \Theta(\phi)=\frac{4\phi(1-\phi)}{W},
  \quad
  \vect n=\frac{\nabla\phi}{|\nabla\phi|+\varepsilon}.
  \label{eq:allen-cahn-target}
\end{equation}
动量部分的连续意图可写成
\begin{equation}
  \partial_t(\rho\vect u^\star)+\nabla\cdot(\rho\vect u^\star\vect u^\star)
  =-\nabla p+\nabla\cdot[\rho\nu(\nabla\vect u^\star+\nabla\vect u^{\star T})]
  +\vect f_g+\vect f_T+\vect f_\sigma,
  \label{eq:momentum-target}
\end{equation}
并配合近似不可压条件 $\nabla\cdot\vect u^\star\approx0$。式~\eqref{eq:allen-cahn-target}--
\eqref{eq:momentum-target} 用于解释模型；后续碰撞、源项和边界公式才是代码实际执行的离散。

\subsection{两个平衡分布}

分布 $f_q$ 携带动态压力和动量，不把 $\sum_qf_q$ 当作密度。令 $p^d=p-p^H$，则
\begin{equation}
  f_q^{\rm eq}
  =3w_qp^d-3p^d\delta_{q0}
  +w_q\rho\left[3\vect c_q\cdot\vect u^\star
  +\frac92(\vect c_q\cdot\vect u^\star)^2
  -\frac32|\vect u^\star|^2\right].
  \label{eq:pressure-equilibrium}
\end{equation}
它满足
\begin{equation}
  \sum_qf_q^{\rm eq}=0,
  \quad \sum_q\vect c_qf_q^{\rm eq}=\rho\vect u^\star,
  \quad \sum_q\vect c_q\vect c_qf_q^{\rm eq}
  =\rho\vect u^\star\vect u^\star+p^d\vect I.
\end{equation}
相场平衡分布为
\begin{equation}
  h_q^{\rm eq}=w_q\phi\left[1+3\vect c_q\cdot\vect u^\star
  +\frac92(\vect c_q\cdot\vect u^\star)^2
  -\frac32|\vect u^\star|^2\right],
  \qquad \sum_qh_q^{\rm eq}=\phi.
  \label{eq:phase-equilibrium}
\end{equation}

\section{初始化：为什么第一步之前已经做了很多工作}

初始化并不是简单把矩形数组填为常数。其顺序为：生成墙、水和冰的初始几何；初始化
$f_q,h_q,\phi$；对相场做 500 步零速预热，使台阶自由液面松弛成宽度约为 $W$ 的弥散界面；
恢复宏观量并做一次全局水量投影；随后才初始化材料冰与世界水热广延量、动态刚体质量性质和
静水参考。

\subsection{冻结的 well-balanced 静水参考}

预热后的相场先在每一水平行上平均，得到参考密度剖面 $\rho_j^H$。以初始自由液面为压力规准
$p^H=0$，采用单元中心梯形积分，例如向下为
\begin{equation}
  p_j^H=p_{j+1}^H-\frac12g_y(\rho_j^H+\rho_{j+1}^H).
  \label{eq:hydrostatic-trapezoid}
\end{equation}
初始化 population 时只写入 $p^d=0$。以后体内重力只保留
$(\rho-\rho^H)\vect g$，而冰 cut-link 的动量交换补回同一 $p^H$ 的牵引。这样静水压力、
体内源项和冰面浮力来自一套离散，而不是在水动力之外再添加经验阿基米德力。

\subsection{初始化后的时间层}

在 $t^0$，$f,h,\phi,p,\vect u$、刚体位姿、材料热状态、世界水热状态和连续覆盖率都对应同一
几何。此后每个 \code{step()} 都把这些量从 $n$ 层推进到 $n+1$ 层。热扩散每八步才更新一次，
但位姿光栅、ALE 和水显热对流每一步都执行。

\section{一个完整 \code{step()} 的总览}

\subsection{与源码逐行对应的算法}

\begin{algorithm}[H]
\caption{一个 LBM 时间步的真实调度顺序}
\label{alg:step}
\small
\begin{algorithmic}[1]
\Require $n$ 层流体 distributions、宏观场、刚体状态、材料冰与世界水热广延量
\Ensure 所有状态推进到 $n+1$ 层，且离散水量、热水 aperture 与待处理融水动量一致
\State \code{\_collide\_velocity()}：计算力、黏度和 $f_q^\star$
\State \code{\_collide\_phase()}：计算 $h_q^\star$
\State \code{\_stream()}：流送/反弹/cut-link，并归约冰体水动力冲量
\State \code{\_update\_streamed\_macroscopic\_fields()}：恢复 $\phi,\vect u$
\State \code{\_update\_pressure()}：恢复 $p=p^H+p^d$
\State \code{\_integrate\_rigid\_ice()}：刚体积分与候选位姿更新
\State \code{\_solve\_contact\_and\_rasterize()}：新位姿下连续材料覆盖率、SDF 与整体防穿墙投影
\State \code{\_update\_solid\_mask()}：更新锐冰掩膜
\State \code{\_refill\_changed\_nodes()}：重建本步启闭的 LBM 节点
\State \code{\_advance\_moving\_thermal\_fast()}：ALE 与成对水体积--显热对流
\State $n_{\rm pending}\gets n_{\rm pending}+1$
\If{$n_{\rm pending}\ge N_{\rm slow}$}
  \State \code{\_advance\_moving\_thermal\_slow()}：导热、换热、融化、融水注入
  \State $n_{\rm pending}\gets0$
\EndIf
\State \code{\_correct\_water\_volume()}：有界 logit 相场水量投影
\State \code{\_finalize\_melt\_momentum\_coupling()}：若有待处理慢热步，闭合线/角动量
\State \code{\_synchronize\_moving\_water\_aperture()}：重排热水广延量并恢复温度
\State \code{steps}$\gets$\code{steps}$+1$
\end{algorithmic}
\end{algorithm}

\subsection{算法框图}

\begin{figure}[H]
\centering
\begin{tikzpicture}[node distance=3.3mm]
  \node[flow] (c) {流体局地更新：$f$ 碰撞 $\rightarrow h$ 碰撞};
  \node[flow,below=of c] (s) {流送与边界：普通流送 / 固定墙反弹 / 移动冰 cut-link；旁路归约 $\vect J_h,L_h$};
  \node[flow,below=of s] (m) {恢复 $\phi,\vect u,p$；刚体积分得到候选位姿};
  \node[flow,below=of m] (g) {连续光栅与 contact 投影 $\rightarrow$ 锐掩膜 $\rightarrow$ 启闭节点 equilibrium refill};
  \node[flow,below=of g] (f) {FAST：位姿/相场联合 aperture 重映射 $\rightarrow$ 成对 $V^w/S^w$ 迎风对流};
  \node[decision,below=5mm of f] (d) {$n_{\rm pending}\ge N_{\rm slow}$？};
  \node[slow,right=19mm of d,text width=5.3cm] (slow) {SLOW：动量快照；冰/水导热；冰--水换热；潜热融化；融水注入；重算 $m,I,\vect\xi_c$；侵蚀光栅与 refill};
  \node[flow,below=7mm of d] (v) {相场水量投影 $\rightarrow$ 条件融水线/角动量闭合};
  \node[flow,below=of v] (a) {最终 aperture 同步与温度/焓诊断恢复；\code{steps}$\leftarrow$\code{steps}$+1$};
  \draw[arrow] (c)--(s);
  \draw[arrow] (s)--(m);
  \draw[arrow] (m)--(g);
  \draw[arrow] (g)--(f);
  \draw[arrow] (f)--(d);
  \draw[arrow] (d)--node[above,font=\small]{是}(slow);
  \draw[arrow] (slow.south) |- (v.east);
  \draw[arrow] (d)--node[right,font=\small]{否}(v);
  \draw[arrow] (v)--(a);
\end{tikzpicture}
\caption{\code{step()} 的数据流。SLOW 分支只在累计步数达到阈值时执行；新温度在下一步
\code{\_collide\_velocity()} 中进入热浮力，因此没有同一步内的隐式反馈。}
\label{fig:step-flow}
\end{figure}

下面各节按照算法~\ref{alg:step} 的顺序解释每一步的输入、公式、输出和设置该步骤的理由。

\section{步骤一：压力--动量分布碰撞}

\stepname{\_collide\_velocity()}
输入为 $f_q,\phi,\vect u,p,p^H,\rho^H,T$、墙/冰掩膜与冰面速度；输出为本步格子加速度
$\vect F$ 和碰撞后分布 $f_q^\star$。

\subsection{相场差分、化学势与表面张力}

遇到墙、锐冰或域外邻点时，相场邻值回退为本点值，等价于让该方向的相场增量为零。
令 $\phi_q^{(1)},\phi_q^{(2)}$ 为沿 $\vect c_q$ 的一格和两格采样，代码使用
\begin{align}
  \nabla_s\phi&=3\sum_qw_q\vect c_q(\phi_q^{(1)}-\phi_P),\\
  \nabla_g\phi&=(1-\alpha)3\sum_qw_q\vect c_q(\phi_q^{(1)}-\phi_P)
  +\alpha\frac32\sum_qw_q\vect c_q(4\phi_q^{(1)}-\phi_q^{(2)}-3\phi_P),\\
  \nabla^2\phi&=6\sum_qw_q(\phi_q^{(1)}-\phi_P),
  \qquad \alpha=\frac13.
  \label{eq:phase-stencils}
\end{align}
$\nabla_g\phi$ 用于化学势力和相场法向；较短的 $\nabla_s\phi$ 用于密度梯度源和压力恢复。
自由能系数、化学势与毛细力密度为
\begin{align}
  \beta&=\frac{12\sigma}{W},\qquad \kappa=\frac32\sigma W,\\
  \mu&=4\beta\bar\phi(\bar\phi-1)(\bar\phi-\tfrac12)-\kappa\nabla^2\phi,\\
  \vect f_\sigma&=\mu\nabla_g\phi.
  \label{eq:chemical-force}
\end{align}

\subsection{well-balanced 重力与温度浮力}

定义只在水侧逐渐开启的门函数与线性 Boussinesq 密度异常
\begin{equation}
  \chi_w(\phi)=\clip(2\bar\phi-1,0,1),
  \qquad \epsilon_T=-\alpha_T(T-T_{\rm ref}).
\end{equation}
总力密度和代码存储的加速度为
\begin{align}
  \vect f_g&=(\rho-\rho^H)\vect g,\\
  \vect f_T&=\chi_w\rho_w\epsilon_T\vect g,\\
  \vect F&=\frac{\vect f_g+\vect f_T+\vect f_\sigma}{\rho}.
  \label{eq:total-lattice-acceleration}
\end{align}
静水参考部分 $\rho^H\vect g$ 不在 bulk 重复施加；其压力牵引会在移动冰 cut-link 的 MEM 中
补回。温度只改变浮力源，不改写 LBM 惯性密度 $\rho(\phi)$，也不改变冰水质量转换密度。

\subsection{人工黏度和两种松弛率}

速度邻值在流体格取流速、冰格取刚体点速度、墙或域外取零。用 D2Q9 差分得到速度梯度，代码的
涡黏度为
\begin{equation}
  \nu_t=c_d^2\sqrt{2\left[(\partial_xu_x)^2+(\partial_yu_y)^2
  +\frac12(\partial_yu_x+\partial_xu_y)^2\right]}.
  \label{eq:eddy-viscosity}
\end{equation}
局部主松弛时间
\begin{equation}
  \nu=\nu_0(\phi)+\nu_t,
  \qquad \tau_p=\frac12+3\nu,
  \qquad \omega_p=\tau_p^{-1}.
\end{equation}
为抑制空气界面和锐冰邻域中未解析的剪切振荡，只对偏应力矩设置下限
\begin{equation}
  \tau_s=\max\left\{\tau_p,
  \frac12+(\tau_{\rm int}-\tfrac12)
  \max(1-\bar\phi,I_{\rm near\,ice})\right\},
  \qquad \omega_s=\tau_s^{-1}.
  \label{eq:shear-floor}
\end{equation}
$I_{\rm near\,ice}$ 表示一格邻域是否存在锐冰；默认 $\tau_{\rm int}=0.8$。

\subsection{源项与原始矩 MRT 碰撞}

用 $\vect u^\star$ 计算式~\eqref{eq:pressure-equilibrium}，离散源为
\begin{equation}
  S_q=3w_q\left[
  \vect c_q\cdot(\rho\vect F)
  +(\rho_w-\rho_a)(\vect c_q\cdot\vect u^\star)
  (\vect c_q\cdot\nabla_s\phi)
  \right].
  \label{eq:momentum-source}
\end{equation}
定义原始单项式矩
\begin{equation}
  M_{ab}=\sum_qc_{qx}^ac_{qy}^bf_q,
  \quad M_{ab}^{\rm eq}=\sum_qc_{qx}^ac_{qy}^bf_q^{\rm eq},
  \quad S_{ab}=\sum_qc_{qx}^ac_{qy}^bS_q.
\end{equation}
代码对 $M_{00},M_{10},M_{01}$ 用 $\omega_p$ 松弛并加梯形源
\begin{equation}
  M_{ab}^\star=M_{ab}-\omega_p(M_{ab}-M_{ab}^{\rm eq})
  +(1-\tfrac12\omega_p)S_{ab}.
  \label{eq:raw-moment-relaxation}
\end{equation}
二阶差 $M_{20}-M_{02}$ 与 $M_{11}$ 改用 $\omega_s$；二阶迹、三阶和四阶 ghost 矩直接置为
$M^{\rm eq}+S/2$。最后以零平移中心反演回 $f_q^\star$。所以这是\textbf{原始矩 MRT}，不是
随局部速度平移的中心矩碰撞。

\section{步骤二：守恒 Allen--Cahn 相场碰撞}

\stepname{\_collide\_phase()}
输入为 $h_q,\phi,\vect u$ 和上一步刚得到的 $\vect F$；输出为 $h_q^\star$。

用式~\eqref{eq:phase-stencils} 的 $\nabla_g\phi$ 定义
\begin{equation}
  \vect n=\frac{\nabla_g\phi}{|\nabla_g\phi|+10^{-14}},
  \qquad \Theta=\frac{4\phi(1-\phi)}{W},
  \qquad R_q=w_q\Theta(\vect c_q\cdot\vect n).
  \label{eq:phase-compression-source}
\end{equation}
相场松弛时间为
\begin{equation}
  \tau_\phi=\frac12+3M,
  \qquad \omega_\phi=\tau_\phi^{-1}.
\end{equation}
与 $f_q$ 不同，$h_q$ 采用以 $\vect u^\star$ 为中心的中心矩
\begin{equation}
  K_{ab}=\sum_q(c_{qx}-u_x^\star)^a(c_{qy}-u_y^\star)^bh_q.
\end{equation}
零阶矩和二阶以上矩直接置为 $K^{\rm eq}+R/2$；两个一阶矩为
\begin{equation}
  K_{10}^\star=K_{10}-\omega_\phi(K_{10}-K_{10}^{\rm eq})
    +(1-\tfrac12\omega_\phi)R_{10},
  \quad
  K_{01}^\star=K_{01}-\omega_\phi(K_{01}-K_{01}^{\rm eq})
    +(1-\tfrac12\omega_\phi)R_{01}.
  \label{eq:phase-central-collision}
\end{equation}
再以同一 $\vect u^\star$ 反演得到 $h_q^\star$。压缩源把弥散界面维持在设定宽度附近，迁移率
$M$ 控制相场扩散；全局水量误差不在此处强行添加源，而在步骤十二的有界投影中闭合。

\section{步骤三：流送、三类边界与冰体动量交换}

\stepname{\_stream()}
输入为 $f_q^\star,h_q^\star$、旧位姿锐冰/SDF、$p^H$；输出为下一层 $f_q,h_q$，并归约一整个
格子步的冰体水动力冲量 $\vect J_h$ 与角冲量 $L_h$。它们已经是冲量，不应再乘 $\Delta t$。

\subsection{三种链路分支}

\begin{figure}[H]
\centering
\begin{tikzpicture}[node distance=7mm and 6mm]
  \node[decision] (q) {目的邻格是什么？};
  \node[branch,below left=of q] (normal) {活动流体格\\直接 push streaming};
  \node[branch,below=of q] (wall) {固定墙或域外\\半格反弹};
  \node[branch,below right=of q] (ice) {移动锐冰格\\cut-link + 移动反弹 + MEM};
  \draw[arrow] (q)--(normal);
  \draw[arrow] (q)--(wall);
  \draw[arrow] (q)--(ice);
\end{tikzpicture}
\caption{\code{\_stream()} 对每个活动单元、每个 D2Q9 方向所作的分支。}
\label{fig:stream-branches}
\end{figure}

若 $P+\vect c_q$ 是活动格，则
\begin{equation}
  f_q(P+\vect c_q,n+1)=f_q^\star(P,n),
  \qquad h_q(P+\vect c_q,n+1)=h_q^\star(P,n).
\end{equation}
若目的地是固定墙或数组外，则作 on-site 半格反弹
\begin{equation}
  f_{\bar q}(P,n+1)=f_q^\star(P,n),
  \qquad h_{\bar q}(P,n+1)=h_q^\star(P,n),
\end{equation}
其中 $\vect c_{\bar q}=-\vect c_q$。

\subsection{移动冰面 cut-link}

令流体端 SDF 为 $d_f$、冰端为 $d_s$，界面在链上的比例和交点为
\begin{equation}
  \eta=\clip\left(\frac{d_f}{d_f-d_s+\varepsilon},0.05,0.95\right),
  \qquad \vect x_b=\vect x_P+\eta\vect c_q.
  \label{eq:cut-location}
\end{equation}
由式~\eqref{eq:body-point-velocity} 得到 $\vect u_b(\vect x_b)$。压力--动量分布在所有移动冰面
cut-link 上统一使用 Tao 等人的 single-node 二阶曲壁条件
~\cite{tao2018}。定义碰撞前非平衡量
\begin{equation}
  f_q^{\rm neq}(P,n)=f_q(P,n)-f_q^{\rm eq}(p^d_P,\rho_P,\vect u_P^\star),
\end{equation}
则反向未知分布为
\begin{equation}
  f_r=\frac{f_{\bar q}^{\rm eq}(p^d_P,\rho_P,\vect u_b)
  +f_q^{\rm neq}(P,n)+\eta f_{\bar q}^\star(P,n)}{1+\eta}.
  \label{eq:single-node-cut-link}
\end{equation}
式中所有分布都来自 $P$，不读取后退点 $P-\vect c_q$；代码在碰撞时缓存
$f_q^{\rm neq}(P,n)$，避免 push streaming 原地覆盖 $f_q(P,n)$。壁速已完整包含在
$f_{\bar q}^{\rm eq}(\vect u_b)$ 中，因此不再叠加额外的线性壁速修正。
空气、自由液面弥散带和水中的冰面链接使用同一公式，不再按 $\phi$ 切换到简单反弹。相场分布
始终采用
\begin{equation}
  h_r=h_q^\star(P)-6w_q\phi\,\vect c_q\cdot\vect u_b,
\end{equation}
不做 single-node 曲壁处理。

\subsection{逐链 MEM 冲量}

静水参考在交点线性插值
\begin{equation}
  p_b^H=(1-\eta)p_P^H+\eta p_{P+\vect c_q}^H.
\end{equation}
该链传给冰体的离散冲量与关于当前质心的角冲量为
\begin{align}
  \vect J_q&=(\vect c_q-\vect u_b)f_q^\star
  -(-\vect c_q-\vect u_b)f_r+6w_qp_b^H\vect c_q,\\
  L_q&=(\vect x_b-\vect X_c)\times\vect J_q.
  \label{eq:mem-link-impulse}
\end{align}
每个边界流体单元先在本地求和，再用 f64 原子归约
$\vect J_h=\sum\vect J_q$、$L_h=\sum L_q$。最后一项补回与冻结静水参考一致的牵引。

\section{步骤四与五：宏观量和压力恢复}

\subsection{流送后宏观场}

\stepname{\_update\_streamed\_macroscopic\_fields()}
对活动格恢复
\begin{equation}
  \phi_P=\sum_qh_{Pq},
  \qquad
  \rho_P=\rho(\phi_P),
  \qquad
  \vect u_P=\frac{1}{\rho_P}\sum_q\vect c_qf_{Pq}.
  \label{eq:macro-recovery}
\end{equation}
此处 $\vect u$ 还没有加半力。冰内显示速度设为当前刚体点速度，墙内设为零；非活动格的
$\vect F$ 清零。

\subsection{动态压力恢复}

\stepname{\_update\_pressure()}
用短梯度 $\nabla_s\rho=(\rho_w-\rho_a)\nabla_s\phi$ 和
$s_0=-\tfrac32w_0|\vect u^\star|^2$，动态压力为
\begin{equation}
  p^d=\frac35\left[
  \sum_{q\ne0}f_q+\frac12\vect u^\star\cdot\nabla_s\rho+\rho s_0
  \right],
  \qquad p=p^H+p^d.
  \label{eq:pressure-recovery}
\end{equation}
压力在刚体移动之前恢复，供快照、下一次碰撞和可能发生的节点 refill 使用。

\section{步骤六：刚体积分、连续光栅与整体防穿墙投影}

\stepname{\_integrate\_rigid\_ice()}
输入包括 cut-link 归约的水动力线冲量 $\vect J_h$、水动力角冲量 $L_h$、当前动态质量性质和旧位姿；
输出包括质心速度 $\vect V$（即符号表中的 $\vect V_c$）、角速度 $\omega$、
质心位置 $\vect X_c$、转角 $\theta$ 和材料参考原点 $\vect X_0$，并消费、清零本步水动力累加器。
本节用积分前时间层上标 $n$ 和积分后时间层上标 $n+1$ 区分状态。质量、惯量、位置、速度和冲量
均使用格子单位，一个刚体格子步的时间增量为 1。

\subsection{侵蚀后的质量、质心和惯量}

以材料单元索引 $\alpha$ 遍历材料网格，记该单元的当前固冰物理质量为 $m_\alpha^s$、参考水密度为
$\rho_w$、物理网格边长为 $\Delta x$，并把换算后的格子质量记为 $\widehat m_\alpha$：
\begin{equation}
  \widehat m_\alpha=\frac{m_\alpha^s}{\rho_w\Delta x^2}.
\end{equation}
该式用单位出平面厚度下一个满水网格单元的参考质量 $\rho_w\Delta x^2$ 归一化固冰物理质量，
从而得到刚体积分核使用的格子质量。

记材料单元中心相对参考原点的格子坐标为 $\vect\xi_\alpha$、材料固相面积率为 $f_\alpha^s$、
刚体总格子质量为 $m_b$、材料坐标质心为 $\vect\xi_c$，并记关于该质心的格子极惯量为 $I_b$。
这些质量性质归约为
\begin{align}
  m_b&=\sum_\alpha\widehat m_\alpha,\\
  \vect\xi_c&=\frac1{m_b}\sum_\alpha\widehat m_\alpha\vect\xi_\alpha,\\
  I_b&=\sum_\alpha\widehat m_\alpha
  \left(|\vect\xi_\alpha-\vect\xi_c|^2+\frac{f_\alpha^s}{6}\right).
  \label{eq:body-mass-moments}
\end{align}
第一式由材料质量的可加性得到零阶矩，第二式用质量一阶矩除以总质量得到质心，第三式由平行轴
定理把各单元的质心偏置惯量与自身极惯量相加。自身极惯量项 $f_\alpha^s/6$ 来自把部分材料单元
近似成面积为 $f_\alpha^s$、边长为 $\sqrt{f_\alpha^s}$ 的等面积正方形。质量只允许单调下降；
当质量或惯量低于数值阈值时，机械刚体退役，但剩余热材料仍可继续存在。

\subsection{基于物理质量和惯量的自由预测}

记格子重力速度增量为 $\vect g$，平移和转动阻尼系数为 $d_v,d_\omega$。墙面限制之前的自由速度预测为
\begin{align}
  \widetilde{\vect V}&=d_v\left[\vect V^n+
  \frac{\vect J_h+m_b^n\vect g}{m_b^n}\right],\\
  \widetilde\omega&=d_\omega\left[\omega^n+
  \frac{L_h}{I_b^n}\right].
  \label{eq:rigid-free-predictor}
\end{align}
水动力和重力在一个格子步内作为冲量相加，因此不再另乘物理时间步。默认
$d_v=1$、$d_\omega=0.999$；式~\eqref{eq:rigid-free-predictor} 的结果直接作为候选
位姿更新的平移速度和角速度。平移速度和角速度不会在刚体积分阶段被人为设置额外上限；随后只在需要时
由墙面整体投影移除指向墙外的法向平移分量。

\subsection{按当前锐界支撑夹住整体位置}

自由预测先给出候选位姿
\begin{equation}
  \vect X_{c,*}^{n+1}=\vect X_c^n+\widetilde{\vect V},
  \qquad
  \theta^{n+1}=\theta^n+\widetilde\omega.
  \label{eq:free-pose-update}
\end{equation}
程序在这个候选位姿上重新构造当前侵蚀冰体的未裁墙锐界：材料固相率按与世界锐掩膜相同的
双线性规则采样，并以同一阈值 $f_{\rm th}=0.5$ 选中世界单元。选中单元按完整单元的四条面计入
支撑范围。该支撑归约不套用墙掩膜，因此即使候选位置已经穿过墙面，也能看到越界部分并把整体
移回容器。

记候选角度下锐界相对质心的四个支撑极值为
$r_x^-,r_x^+,r_y^-,r_y^+$。若固定墙厚为 $b=\code{boundary\_cells}$，则物理墙与活动域的内边界
位于 $x=b$、$x=N_x-b$、$y=b$ 和 $y=N_y-b$；质心允许区间为
\begin{align}
  x_{\min}&=b-r_x^-,
  &x_{\max}&=N_x-b-r_x^+,\\
  y_{\min}&=b-r_y^-,
  &y_{\max}&=N_y-b-r_y^+.
  \label{eq:body-center-wall-bounds}
\end{align}
最终质心逐坐标夹到这个区间：
\begin{equation}
  X_{c,x}^{n+1}=\clip(X_{c,*,x}^{n+1},x_{\min},x_{\max}),
  \qquad
  X_{c,y}^{n+1}=\clip(X_{c,*,y}^{n+1},y_{\min},y_{\max}).
  \label{eq:body-position-clamp}
\end{equation}
因此约束对象是当前剩余冰体的整体锐界支撑，且边界就是物理墙面；没有向流体域内额外保留两格或
其他 clearance。若当前已不存在可归约的锐界支撑，位置投影为恒等操作。

\subsection{只截断指向墙外的平移速度分量}

墙面速度处理检查候选质心是否达到或越过式~\eqref{eq:body-center-wall-bounds}。左右墙仅修改
$x$ 分量：
\begin{equation}
  V_x^{n+1}=
  \begin{cases}
    \max(\widetilde V_x,0),&X_{c,*,x}^{n+1}\le x_{\min},\\
    \min(\widetilde V_x,0),&X_{c,*,x}^{n+1}\ge x_{\max},\\
    \widetilde V_x,&\text{其他情况},
  \end{cases}
  \label{eq:wall-normal-x-clamp}
\end{equation}
下、上墙仅修改 $y$ 分量：
\begin{equation}
  V_y^{n+1}=
  \begin{cases}
    \max(\widetilde V_y,0),&X_{c,*,y}^{n+1}\le y_{\min},\\
    \min(\widetilde V_y,0),&X_{c,*,y}^{n+1}\ge y_{\max},\\
    \widetilde V_y,&\text{其他情况}.
  \end{cases}
  \label{eq:wall-normal-y-clamp}
\end{equation}
于是撞左墙只删除向左的 $V_x<0$，撞右墙只删除向右的 $V_x>0$；撞底墙或顶墙同理只删除相应
向外的 $V_y$。切向平移速度保持不变，角速度也完全不由墙面处理修改：
\begin{equation}
  \omega^{n+1}=\widetilde\omega.
\end{equation}
若候选位姿同时越过两面墙，两坐标规则独立应用。该边界只是防止整体穿墙的运动学保护，不计算
接触点冲量、摩擦、恢复或角冲量，也没有 PGS 迭代、未分辨 wall-film 或任何近壁特殊处理。

最后用保留的候选角度和投影后的质心恢复材料参考原点：
\begin{equation}
  \vect X_0^{n+1}=\vect X_c^{n+1}
  -\vect R(\theta^{n+1})\vect\xi_c.
\end{equation}
该式保证投影平移整个当前冰体，而不改变其相对材料几何。

\section{步骤七至九：新位姿光栅、锐界与 equilibrium refill}

\subsection{连续材料覆盖与 SDF}

\stepname{\_solve\_contact\_and\_rasterize()}
对世界单元中心 $\vect x_P$ 做逆刚体变换
\begin{equation}
  \vect\xi_P=\vect R(\theta)^T(\vect x_P-\vect X_0),
  \qquad
  C_P=\sum_{\alpha\in\mathcal B(P)}\omega_{\alpha P}f_\alpha^s,
  \label{eq:world-raster}
\end{equation}
其中 $\mathcal B(P)$ 是双线性插值的最多四个材料单元。调用前的 $C_P$ 保存为 $C_P^{\rm old}$，
新值保存为 $C_P^{\rm new}$，用于锐界判定和热接触；FAST 根据新开口联合重映射水体积和显热。材料矩形内部的 SDF 为
\begin{equation}
  d_P=(f_{\rm th}-C_P)\Delta x,
  \qquad f_{\rm th}=0.5,
  \label{eq:world-sdf}
\end{equation}
矩形外则使用旋转矩形的有向距离。这里 $d_P$ 的单位是米，并直接写入程序唯一的
\code{sdf} 场；cut-link 只使用距离比，因此单位相消。

\subsection{锐冰掩膜更新}

\stepname{\_update\_solid\_mask()}
先保存 $\chi_P^{s,n}$ 到 \code{solid\_prev}，再定义
\begin{equation}
  \chi_P^{s,n+1}=
  \begin{cases}
  1,&\text{刚体仍活动、$P$ 不是墙且 }d_P\le0,\\
  0,&\text{其他情况}.
  \end{cases}
  \label{eq:sharp-mask}
\end{equation}
因此热接触可以在任意 $C_P>0$ 时继续存在，而 LBM 锐冰只在 $C_P\ge0.5$ 时占据节点。

\subsection{启闭节点重建}

\stepname{\_refill\_changed\_nodes()}
若节点由流体变为冰，程序保存裁剪后的 $\phi$ 和动态压力残量 $p-p^H$ 作为覆盖 reservoir，
显示速度改为刚体点速度，$\vect F=0$；其旧 population 因节点已非活动而不参与求解。

若节点由冰变为流体，只从 $3\times3$ 邻域中“本次更新前已经是流体”的八邻格平均
$\phi,p-p^H,\vect u$。这一限制避免同一 GPU kernel 中新节点互相读取造成调度依赖。若无合法旧
流体邻格，则回退到覆盖 reservoir、动态压力储备和当前刚体点速度。然后完整重建
\begin{equation}
  f_q=f_q^{\rm eq}(p^d,\rho(\phi),\vect u),
  \qquad h_q=h_q^{\rm eq}(\phi,\vect u),
  \qquad \vect F=\vect0.
  \label{eq:equilibrium-refill}
\end{equation}
这是一种确定性的 equilibrium refill：它不外推非平衡应力，因为新开放格没有可信的旧动力学历史。

\section{步骤十：每个 LBM 步执行的 FAST 热更新}

\stepname{\_advance\_moving\_thermal\_fast()}
FAST 的职责只有三件事：检查热对流所读速度、对刚体位姿变化做守恒 ALE、对流液水体积和液水
显热。它不做导热或融化。

\subsection{半力速度与 Courant 子步}

热对流速度取活动格的 $\vect u^\star$，其他格为零。调试开关
\code{\small\_thermal\_stability\_check\_enabled} 置为真时，程序拒绝 NaN/Inf，并要求全活动域
$|u_x^\star|+|u_y^\star|\le1$；水侧 CFL 最大值只在 $\phi>\varepsilon_\phi$ 的合法格统计。
该 LBM 状态扫描默认关闭，但对流所需的速度场和水侧 CFL 归约仍然每步执行。
物理对流子步上限与子步数为
\begin{equation}
  \Delta t_{\rm adv,max}
  =Co_{\max}\frac{\Delta t_{\rm LBM}}
  {\max_P(|u_x^\star|+|u_y^\star|)},
  \qquad
  N_{\rm adv}=\max\left(1,
  \left\lceil\frac{\Delta t_{\rm LBM}}{\Delta t_{\rm adv,max}}\right\rceil\right).
  \label{eq:advection-substeps}
\end{equation}
这里格子速度乘 $\Delta x/\Delta t_{\rm LBM}$ 后才是物理速度。

\subsection{位姿与相场开口联合重映射}

FAST 直接按第~\ref{sec:aperture} 节的新开口目标重映射 $V^w,S^w$。
不再先按连续覆盖率扫掠一次、再按锐掩膜同步一次。原来的全域供体池虽然闭合总量，
却把冰旁冷水混入整条上升自由面，又把远处热水混入冰旁 refill；这会产生初始水面附近的
人工温度分层。现实现使用局部比显热重构和有界总能量校正，以减少这种空间误差。

\subsection{成对液水面积--显热迎风对流}

联合 aperture 重映射保证上游状态合法。对相邻合法水格
$L,R$ 的面，取物理法向速度、上游含水率和两种通量
\begin{align}
  u_f&=\frac12(\vect u_L^\star+\vect u_R^\star)\cdot\vect n_f
  \frac{\Delta x}{\Delta t_{\rm LBM}},\\
  a_U&=\clip\left(\frac{V_U^w}{\cellarea},0,1\right),\\
  F_f^V&=u_fa_U\Delta x,\\
  F_f^S&=F_f^V\frac{S_U^w}{\max(V_U^w,\varepsilon)}.
  \label{eq:paired-upwind-flux}
\end{align}
$U$ 是由 $\sgn(u_f)$ 决定的上游格。对任意子步 $\delta t$，有限体积更新为
\begin{equation}
  V_P^{w,+}=V_P^w-\delta t\sum_{f\in\partial P}F_f^V,
  \qquad
  S_P^{w,+}=S_P^w-\delta t\sum_{f\in\partial P}F_f^S.
  \label{eq:paired-upwind-update}
\end{equation}
只有面两侧均非墙、非锐冰且 $\phi>\varepsilon_\phi$ 时通量开放。$V^w$ 与 $S^w$ 共用体积通量，
所以均匀温度不会因自由表面小体积分数而凭空产生尖峰。若还有下一个 Courant 子步，立即再做
aperture 同步；最后一次同步故意推迟到相场水量投影之后。

\section{步骤十一：条件 SLOW 导热、换热、融化与融水注入}

\stepname{\_advance\_moving\_thermal\_slow()}
当累计 FAST 步数达到 $N_{\rm slow}=8$ 时执行。总物理时间为
$\Delta t_{\rm S}=N_{\rm slow}\Delta t_{\rm LBM}$。进入前先同步 aperture，并记录慢热前活动流体
线/角动量及累计刚体熔失动量基线，供步骤十三闭合。

\subsection{Fourier 子步与温度恢复}

取最大热扩散率
\begin{equation}
  \alpha_{\max}=\max\left(\frac{k_w}{\rho_wc_{p,w}},
  \frac{k_i}{\rho_ic_{p,i}}\right),
\end{equation}
显式扩散上限和子步数为
\begin{equation}
  \Delta t_{\rm diff,max}=Fo_{\max}\frac{\Delta x^2}{\alpha_{\max}},
  \qquad
  N_{\rm diff}=\max\left(1,
  \left\lceil\frac{\Delta t_{\rm S}}{\Delta t_{\rm diff,max}}\right\rceil\right).
  \label{eq:diffusion-substeps}
\end{equation}
每个扩散子步先由式~\eqref{eq:body-thermal-state} 和~\eqref{eq:water-extensive-state} 恢复温度。

\subsection{冰内、水内与容器壁导热}

材料冰相邻单元面 $f=L|R$ 的热功率为
\begin{equation}
  Q_{f}^{i}=-k_i\min(f_L^s,f_R^s)(T_R-T_L).
  \label{eq:ice-conduction-flux}
\end{equation}
世界水相邻格的开口和热功率为
\begin{equation}
  a_f^w=\min\left(1,\frac{V_L^w}{\cellarea},\frac{V_R^w}{\cellarea}\right),
  \qquad Q_f^w=-k_wa_f^w(T_R-T_L).
  \label{eq:water-conduction-flux}
\end{equation}
二维正方网格的面长 $\Delta x$ 与中心距 $\Delta x$ 相消，所以式中没有显式 $\Delta x$。
显热按 $E\leftarrow E-\delta t\,\mathrm{div}\,Q$ 更新。Dirichlet 壁距离单元中心只有半格，故
入水热功率为
\begin{equation}
  Q_b=2k_wa_P^w(T_b-T_P),
  \qquad a_P^w=\clip(V_P^w/\cellarea,0,1).
  \label{eq:wall-heat-flux}
\end{equation}
绝热壁功率为零。当前默认四壁均绝热；$90\,^{\circ}\mathrm C$ 仅为初始水温。空气没有
独立能量方程，水--气、冰--气和冰--壁直接接触均绝热。

\subsection{离散冰--水界面换热}

谐均导热系数为
\begin{equation}
  k_{iw}=\frac{2k_ik_w}{k_i+k_w}.
\end{equation}
对每个合法水格 $P$，程序检查自身和四个轴向世界位置
$\mathcal N_0(P)=\{P,P\pm\vect e_x,P\pm\vect e_y\}$。若位置 $Q$ 的连续冰覆盖 $C_Q>0$，定义
\begin{equation}
  a_P^w=\min(1,V_P^w/\cellarea),
  \quad a_Q^i=\min(1,C_Q/f_{\rm th}),
  \quad b_{\alpha Q}=\frac{\omega_{\alpha Q}f_\alpha^s}{C_Q}.
\end{equation}
对材料贡献格 $\alpha$ 的请求热量为
\begin{equation}
  \Delta Q_{P,Q,\alpha}^{\rm req}
  =\delta t\,k_{iw}\min(a_P^w,a_Q^i)b_{\alpha Q}
  \max(T_{w,P}-T_{i,\alpha},0).
  \label{eq:interface-heat-request}
\end{equation}
同一水格的全部请求按
\begin{equation}
  \gamma_P=\min\left(1,
  \frac{\max(S_P^w,0)}{\sum_{Q,\alpha}\Delta Q_{P,Q,\alpha}^{\rm req}}
  \right)
\end{equation}
同比缩放；最终同一数值从 $S_P^w$ 扣除并加到 $E_\alpha^s$。因此界面交换在离散上成对，且不会
把水显热降到熔点以下。当前模型只允许较热水向不高于熔点的原始冰传热，不是可逆凝固模型。

\subsection{显热--潜热分配}

令材料格本子步净得热为 $Q_\alpha$。若 $Q_\alpha\le0$ 且仍有冰，则直接
$E_\alpha^s\leftarrow E_\alpha^s+Q_\alpha$。若 $Q_\alpha>0$，先填平固冰显热亏空
\begin{equation}
  Q_{\rm sen}=\min(Q_\alpha,\max[-E_\alpha^s,0]),
  \qquad E_\alpha^s\leftarrow E_\alpha^s+Q_{\rm sen}.
\end{equation}
余热 $R=Q_\alpha-Q_{\rm sen}$ 支付潜热：
\begin{align}
  \Delta m_\alpha&=\min\left(m_\alpha^s,\frac{R}{L}\right),\\
  m_\alpha^s&\leftarrow m_\alpha^s-\Delta m_\alpha,\\
  \Delta V_\alpha^w&=\frac{\Delta m_\alpha}{\rho_w},\\
  \Delta S_\alpha^w&=\max(0,R-L\Delta m_\alpha).
  \label{eq:melt-update}
\end{align}
融水在熔点生成，通常 $\Delta S_\alpha^w=0$；只有材料格完全融尽后仍有余热，余热才随融水进入
液水显热。正残量不会被容差静默删除。

\subsection{融水注入与慢热结束后的几何更新}

融水优先注入本次与该材料格交换热量的合法水格；否则在材料点当前世界位置附近的
$5\times5$ 冻结候选中选最近合法湿格。若仍失败，则先按全域空余容量
$(\cellarea-V_P^w)_+$，再退化按已有 $V_P^w$ 比例守恒分配。体积、显热和质量示踪使用相同权重。

每个热子步末执行 aperture 同步；若还有下一子步，则在同一固定刚体位姿重新光栅化连续侵蚀体。
整个 SLOW 结束后，父求解器重新归约式~\eqref{eq:body-mass-moments}。非对称融化导致材料质心
变化时，剩余刚体采取“无喷射反冲”速度
\begin{equation}
  \vect V_c^+=\vect V_c^-+\omega\vect e_z\times
  \vect R(\theta)(\vect\xi_c^+-\vect\xi_c^-),
  \label{eq:no-ejection-recoil}
\end{equation}
并从步前/步后实际质量矩记录熔失线/角动量。之后更新连续覆盖、锐冰、侵蚀造成的 refill 和新的
相场水量目标，并把“融水动量待闭合”标志设为真。

\section{步骤十二：有界相场水量投影}

\stepname{\_correct\_water\_volume()}
这个步骤修改 LBM 相场总水量；它与第~\ref{sec:aperture} 节只重排热水 $V^w,S^w$ 的 aperture
同步是两个不同约束。

\subsection{由真实冰水密度得到目标}

当前剩余固冰的格子体积为
\begin{equation}
  V_i(t)=\frac{m_b(t)}{\rho_i/\rho_w}.
\end{equation}
相场目标水量（以世界单元面积为单位）为
\begin{equation}
  V_w^\star(t)=V_w^0-\frac{\rho_i}{\rho_w}[V_i(t)-V_i(0)]
  =V_w^0+\frac{M_s(0)-M_s(t)}{\rho_w\Delta x^2}.
  \label{eq:water-volume-target}
\end{equation}
它只随材料质量变化，不随刚体平移或旋转引起的锐掩膜格数抖动。由于
$\rho_i/\rho_w=0.917$，融化相同冰体积只生成较小液水体积，差额由自由液面移动吸收。

\subsection{尾部规范化与真实界面活动集}

先将 $\phi$ 裁到 $[0,1]$；默认 $\varepsilon_\phi=10^{-3}$，故
\begin{equation}
  \phi\le\varepsilon_\phi\Rightarrow\phi=0,
  \qquad
  \phi\ge1-\varepsilon_\phi\Rightarrow\phi=1.
  \label{eq:phase-canonicalization}
\end{equation}
改变 $\phi$ 时，以新旧平衡态之差提升 $f_q,h_q$，从而保持动态压力、速度及健康的非平衡部分：
\begin{align}
 f_q'&=f_q+f_q^{\rm eq}(p^d,\rho(\phi'),\vect u)
             -f_q^{\rm eq}(p^d,\rho(\phi),\vect u),\\
 h_q'&=h_q+h_q^{\rm eq}(\phi',\vect u)
             -h_q^{\rm eq}(\phi,\vect u).
 \label{eq:phase-population-lift}
\end{align}
对 $h_0$ 另补舍入余量，使 $\sum_qh_q'=\phi'$。严格 bulk 中若 $\sum_q|h_q|>2$，则直接按
$\vect u^\star$ 重建 $h_q^{\rm eq}$，消除以大数相消表示零阶矩的病态 ghost 模式。

若当前水量已满足相对容差，则结束；否则，只把与局部 $\phi=1/2$ 交叉相连的弥散格作为种子，
再在 $\varepsilon_\phi<\phi<1-\varepsilon_\phi$ 内按固定半径作八邻域膨胀，得到可调活动集
$\mathcal I$。这样不会在远处 bulk 空气中均匀加水。

\subsection{logit 位移与受保护求根}

从优化角度，该步骤是在可调集合上求 Bernoulli 相对熵最小的投影
\begin{equation}
 \min_{\phi'}\sum_{P\in\mathcal I}
 \left[\phi_P'\ln\frac{\phi_P'}{\phi_P}
 +(1-\phi_P')\ln\frac{1-\phi_P'}{1-\phi_P}\right],
 \qquad
 \sum_{P\in\mathcal A}\phi_P'=V_w^\star .
\end{equation}
其解是共同 logit 位移
\begin{equation}
  \phi_P(\lambda)=\frac{\phi_Pe^\lambda}
  {1-\phi_P+\phi_Pe^\lambda},
  \qquad
  \frac{\dd\phi_P}{\dd\lambda}=\phi_P(\lambda)[1-\phi_P(\lambda)].
  \label{eq:logit-map}
\end{equation}
$\phi=0,1$ 是不动点。根方程为
\begin{equation}
  G(\lambda)=\sum_{P\in\mathcal A\setminus\mathcal I}\phi_P
  +\sum_{P\in\mathcal I}\phi_P(\lambda)-V_w^\star=0.
  \label{eq:water-root}
\end{equation}
搜索区间由允许界面平移限制
\begin{equation}
  |\lambda|\le\frac{4\,\Delta s_{\max}}{W},
  \qquad \Delta s_{\max}=0.5\ \text{格}.
\end{equation}
代码先把候选值量化成实际 f32 $\phi$ 并规范化端点，再以 f64 归约；Newton 步若离开括区间则
退回二分，最多 64 次。应用同一映射和式~\eqref{eq:phase-population-lift} 后，最多再做八轮
\begin{equation}
  \delta\phi_P=R\frac{\phi_P(1-\phi_P)}
  {\sum_{Q\in\mathcal F}\phi_Q(1-\phi_Q)}
  \label{eq:phase-residual-correction}
\end{equation}
消除 f32 端点造成的小余量。没有真实界面、目标超容量或所需平移过大时显式报错，而不污染 bulk。

\section{步骤十三：熔化引起的线动量与角动量闭合}

\stepname{\_finalize\_melt\_momentum\_coupling()}
只有 SLOW 刚执行且待处理标志为真时才工作。它不是指定一个局部“融水喷射速度”，而是把慢热、
侵蚀 refill 和相场投影合在一起造成的全局一阶矩变化闭合到刚体实际熔失的离散动量。

慢热前活动流体的线/角动量记为 $(\vect P_f^-,L_f^-)$；相场投影后的暂态值为
$(\vect P_f^c,L_f^c)$。由刚体质量矩步前/步后差得到目标熔失量 $(\Delta\vect P_b,\Delta L_b)$，
需要补入流体的缺口为
\begin{align}
  \Delta\vect P&=\Delta\vect P_b-(\vect P_f^c-\vect P_f^-),\\
  \Delta L&=\Delta L_b-(L_f^c-L_f^-).
  \label{eq:melt-momentum-gap}
\end{align}
在所有合法湿格取权 $a_P=\clip(\phi_P,0,1)$，定义
\begin{equation}
  A=\sum_Pa_P,
  \qquad \bar{\vect x}=\frac1A\sum_Pa_P\vect x_P,
  \qquad J_A=\sum_Pa_P|\vect x_P-\bar{\vect x}|^2.
\end{equation}
同时满足两个线动量矩和一个角动量矩的最小结构修正为
\begin{align}
  \delta\vect p_P&=a_P\left[
  \frac{\Delta\vect P}{A}
  +\gamma\vect e_z\times(\vect x_P-\bar{\vect x})\right],\\
  \gamma&=\frac{\Delta L-\bar{\vect x}\times\Delta\vect P}{J_A}.
  \label{eq:three-moment-lift}
\end{align}
它满足
\begin{equation}
  \sum_P\delta\vect p_P=\Delta\vect P,
  \qquad
  \sum_P\vect x_P\times\delta\vect p_P=\Delta L.
\end{equation}
每格通过 D2Q9 lift 写入
\begin{equation}
  \delta f_{Pq}=3w_q\vect c_q\cdot\delta\vect p_P,
  \qquad
  \vect u_P\leftarrow\vect u_P+\frac{\delta\vect p_P}{\rho_P}.
  \label{eq:d2q9-momentum-lift}
\end{equation}
由 D2Q9 各向同性，$\sum_q\delta f_{Pq}=0$ 且
$\sum_q\vect c_q\delta f_{Pq}=\delta\vect p_P$，所以不改变相场或 $f$ 的零阶矩。最后用两个相异
湿格构造力--力偶，做两轮 f32 舍入余量消除；随后更新累计诊断并清除待处理标志。

\section{步骤十四：最终热水 aperture 同步与温度恢复}
\label{sec:aperture}

\stepname{\_synchronize\_moving\_water\_aperture()}
以世界单元索引 $P$ 遍历固定网格，并以离散速度方向索引 $q$ 区分 D2Q9 方向。该步骤只重排
热模块中世界单元的液水面积 $V_P^w$ 和水显热广延量 $S_P^w$，不修改 LBM 的水相场 $\phi_P$、
压力--动量分布 $f_{Pq}$ 或相场分布 $h_{Pq}$。

记墙掩膜为 $\chi_P^w$、锐冰掩膜为 $\chi_P^s$、水相判定阈值为 $\varepsilon_\phi$；本步骤使用的
水相场值 $\phi_P$ 已在步骤十二裁到区间 $[0,1]$，实现中还会再次防御性裁剪。把能够容纳热水的
单元组成合法水侧集合 $\mathcal W$：
\begin{equation}
  \mathcal W=\{P\mid \chi_P^w=0,\ \chi_P^s=0,
  \ \phi_P>\varepsilon_\phi\}.
\end{equation}
该集合式来自三个可写入条件的交集：世界单元既不能是墙或锐冰，又必须含有超过判定阈值的水相。

记世界单元边长为 $\Delta x$，相应的单元满面积为 $\cellarea=\Delta x^2$；将所有世界单元的
液水面积之和记为总热水面积 $V_\Sigma$，将合法集合内的相场加权面积记为相场基准容量 $B$，
将合法水侧单元数 $|\mathcal W|$ 对应的满格面积记为满格容量 $C$：
\begin{equation}
  V_\Sigma=\sum_PV_P^w,
  \qquad B=\cellarea\sum_{P\in\mathcal W}\phi_P,
  \qquad C=\cellarea|\mathcal W|.
\end{equation}
这三个归约式分别来自对现有热水面积求和、把相场水分数乘以单元面积后求和，以及把每个合法单元
都按满格计算容量。

把合法水侧单元应达到的液水面积记为目标面积 $D_P$，非法水侧单元的目标面积取零；当总热水
面积 $V_\Sigma$ 超过相场基准容量 $B$ 时，再用剩余开口填充比例 $s$ 分配超出的水。满足总面积
守恒和逐格面积有界的目标仅在总热水面积不超过满格容量，即 $0\le V_\Sigma\le C$ 时存在，其形式为
\begin{equation}
  D_P=
  \begin{cases}
  \cellarea\dfrac{V_\Sigma}{B}\phi_P,
    &P\in\mathcal W,\ V_\Sigma\le B,\\[2mm]
  \cellarea\left[\phi_P+s(1-\phi_P)\right],
    &P\in\mathcal W,\ V_\Sigma>B,\\
  0,&P\notin\mathcal W,
  \end{cases}
  \qquad s=\frac{V_\Sigma-B}{C-B},
  \label{eq:aperture-target}
\end{equation}
第一支把相场基准容量等比例缩小到总热水面积，第二支先保留全部相场基准容量，再把超出部分按
各格剩余开口等比例填入，第三支禁止向非法水侧单元放水；要求目标面积之和等于总热水面积便
得到比例系数 $s$。
实现把退化分母按零贡献处理，并把目标面积 $D_P$ 裁到单元允许区间 $[0,\cellarea]$；若总热水
面积 $V_\Sigma$ 超过满格容量 $C$，则报告合法水侧容量不足。

已有水的格保留局部比显热 $e_P=S_P^w/V_P^w$。新开放格从最近的湿邻域外推，最多搜索三圈，
权重为邻格水面积除以距离平方；没有局部水邻居的孤立新组件才回退到当前全水域平均比显热。
这一步只读取冻结的旧水状态，不使用初始水温或空气温度作为能量来源。令
$\widehat S_P=D_P e_P$，并取旧水比显热的全域极值 $e_-,e_+$，定义
\begin{equation}
 \delta E=\sum_P S_P^w-\sum_P\widehat S_P,\quad
 C_-=\sum_P(\widehat S_P-D_Pe_-),\quad
 C_+=\sum_P(D_Pe_+-\widehat S_P).
\end{equation}
最后令 $V_P^{w,+}=D_P$，并作一次有界的总显热校正：
\begin{equation}
 S_P^{w,+}=\widehat S_P+
 \begin{cases}
 \dfrac{\delta E}{C_-}(\widehat S_P-D_Pe_-),&\delta E<0,\\[2mm]
 \dfrac{\delta E}{C_+}(D_Pe_+-\widehat S_P),&\delta E\ge0.
 \end{cases}
\end{equation}
零容量不参与校正，浮点实现把系数限制在 $[-1,0]$ 或 $[0,1]$。
记同步前、后的液水面积分别为 $V_P^{w,-}$ 和 $V_P^{w,+}$，同步前、后的水显热广延量分别为
$S_P^{w,-}$ 和 $S_P^{w,+}$，则有
\begin{equation}
  \sum_PV_P^{w,+}=\sum_PV_P^{w,-},
  \qquad
  \sum_PS_P^{w,+}=\sum_PS_P^{w,-},
  \qquad 0\le V_P^{w,+}\le\cellarea.
  \label{eq:aperture-conservation}
\end{equation}
前两条等式分别来自有界水量目标及总显热校正；这些关系在精确算术下成立，实际浮点误差
另存为守恒残量诊断。该算法还保持均匀水温，并且不产生新的温度极值。全局校正仍然是非局部
数值热再分配，不能视为逐面精确的几何输运；它用有限的局部精度换取每步不必迭代解几何通量。
累计校正绝对量由 \code{aperture\_energy\_correction\_abs\_j\_m} 单独记录。

最后，记水密度为 $\rho_w$、水定压比热为 $c_{p,w}$、熔点温度为 $T_m$、同步后的液水温度为
$T_{w,P}^{+}$，并记温度恢复所用的面积容差为 $V_{\rm tol}=10^{-30}\,\mathrm m^2$。把液水显热
定义式~\eqref{eq:water-extensive-state} 对液水温度反解，得到
\begin{equation}
  T_{w,P}^{+}=T_m+\frac{S_P^{w,+}}{\rho_wc_{p,w}V_P^{w,+}},
  \qquad V_P^{w,+}>V_{\rm tol}.
  \label{eq:water-temperature-recovery}
\end{equation}
该式只是由同步后仍守恒的液水面积与水显热恢复强度量温度，并未再次增加、移除或输运热量。
对于不超过面积容差的同步后液水面积 $V_P^{w,+}\le V_{\rm tol}$，程序把水温恢复字段回退为初始
空气温度；随后合成世界显示温度时，材料冰或墙还会按各自的显示优先级覆盖该回退值。最后一并
合成世界液相率和式~\eqref{eq:diagnostic-enthalpy} 定义的诊断体积焓 $H_P$；其中世界显示温度在
下一 LBM 步的步骤一才进入浮力。

\section{多速率时序：FAST 与 SLOW 为什么不能互换}

\begin{figure}[H]
\centering
\begin{tikzpicture}[x=1.28cm,y=0.75cm,font=\small]
  \foreach \k in {1,...,8}{
    \draw[fill=waterblue] (\k,1) rectangle +(0.9,0.55);
    \node at (\k+0.45,1.275) {F\k};
  }
  \draw[fill=heatred] (8,0) rectangle +(0.9,0.55);
  \node at (8.45,0.275) {S};
  \draw[arrow] (1,1.8)--(8.9,1.8);
  \node[anchor=west] at (1,2.1) {LBM 位姿、ALE 与液水显热对流：每步};
  \node[anchor=west] at (1,-0.35) {导热、界面换热与融化：累计 8 步后一次};
  \draw[dashed] (8.9,0.55)--(8.9,1);
\end{tikzpicture}
\caption{默认 $N_{\rm slow}=8$ 的多速率时序。F 表示 FAST，S 表示 SLOW。}
\label{fig:multirate}
\end{figure}

刚体每一步都会移动，所以流体开口与温度重映射每步执行；否则对流会读到已经关闭的水单元。
导热和相变变化较慢，可累计物理时间后按 Fourier 条件再分子步。水显热对流仍必须
每个 LBM 步跟随瞬时速度执行。

\section{守恒诊断与数值验收}

\subsection{质量与热能}

热系统总质量为
\begin{equation}
  M_T(t)=\sum_\alpha m_\alpha^s(t)+\rho_w\sum_PV_P^w(t),
  \qquad R_M=M_T(t)-M_T(0).
  \label{eq:total-mass}
\end{equation}
以熔点固冰为能量零点，总热能和边界闭合残差为
\begin{align}
  E_T(t)&=\sum_\alpha E_\alpha^s(t)+\sum_PS_P^w(t)
  +L[M_s(0)-M_s(t)],\\
  R_E(t)&=E_T(t)-E_T(0)-Q_b(t).
  \label{eq:total-energy}
\end{align}
这里不含流体/刚体动能，也不包含机械能转热。
四壁绝热时 $Q_b=0$。若有足够热量完全融化，最终混合温度应为
\begin{equation}
 T_{\rm eq}=T_m+
 \frac{m_{w,0}c_{p,w}(T_{w,0}-T_m)+m_{i,0}c_{p,i}(T_{i,0}-T_m)-m_{i,0}L}
 {(m_{w,0}+m_{i,0})c_{p,w}}.
\end{equation}
默认实际水面积扣除墙层后为 $10998\Delta x^2$，故 $T_{\rm eq}\simeq76.64\,^{\circ}\mathrm C$。
完全融化增加的水面积为 $0.917\times32^2\Delta x^2=5.8688\times10^{-5}\,\mathrm m^2$。
若式中分子为负，则热量不足以完全融化；不能把该值当作全液态终温。

\subsection{各层投影的独立残差}

程序分别记录而不是混成一个“总误差”：
\begin{align}
  R_\phi&=\sum_{P\in\mathcal A}\phi_P-V_w^\star,\\
  R_V^{\rm ALE}&=R_V^{\rm ap},
  &R_S^{\rm ALE}&=R_S^{\rm ap},\\
  R_V^{\rm ap}&=\sum_PV_P^{w,+}-\sum_PV_P^{w,-},
  &R_S^{\rm ap}&=\sum_PS_P^{w,+}-\sum_PS_P^{w,-},\\
  R_m^{\rm melt}&=\Delta M_{w,\rm inject}-\Delta M_{i,\rm melt}.
  \label{eq:projection-residuals}
\end{align}
ALE 名称为历史输出兼容别名，现与联合 aperture 同步残差相同。
还记录满格容量余量 $C-\sum_PV_P^w$。累计融水动量残差为
\begin{equation}
  \vect R_P=\sum_k\Delta\vect P_{f,k}-\sum_k\Delta\vect P_{b,k},
  \qquad
  R_L=\sum_k\Delta L_{f,k}-\sum_k\Delta L_{b,k}.
  \label{eq:melt-momentum-residuals}
\end{equation}
\code{synchronize\_diagnostics()} 只在快照边界做输出归约，避免每个 LBM 步为了写日志付出全域归约成本。

\subsection{建议的排错顺序}

如果结果突然异常，建议先判断是哪一层约束失效：
\begin{enumerate}
  \item 若打开 \code{\small\_thermal\_stability\_check\_enabled}，先查 NaN/Inf 和
  $|u_x|+|u_y|$，排除 LBM 已离开稳定包络；
  \item 查 $R_\phi$，判断自由液面水量投影是否可达；
  \item 检查新开放格温度与邻域是否相符，并检查累计显热校正量，不能仅以总量闭合验收温度场；
  \item 查 $R_V^{\rm ap},R_S^{\rm ap}$ 与容量余量，判断热水是否有合法开口；
  \item 查 $R_m^{\rm melt}$，判断潜热生成与融水注入是否配对；
  \item 最后查 $\vect R_P,R_L$，判断侵蚀 refill、相场修正与融水动量 lift 的三矩闭合。
\end{enumerate}
冰体到墙时应同时检查投影前后质心、式~\eqref{eq:body-center-wall-bounds} 的安全区间以及被截断的
法向速度分量。当前墙面操作不会生成接触冲量，也不会改切向速度或角速度；若出现
向内“横向爆冲”或角速度突变，应继续检查 cut-link 水动力、LBM 稳定性和融水动量闭合，而不能归因于
不存在的墙面冲量求解器。

\section{稳定性、精度和模型边界}

\begin{itemize}
  \item LBM 的低马赫条件没有由热子循环修复；若 $|\vect u|$ 过大，应该重新选择格子尺度或
  参考速度，而不是无限增加热对流子步。
  \item $\tau_p>1/2$ 是正运动黏度的基本条件；界面剪切 floor 只提高部分应力矩的耗散。
  \item FAST 的一阶迎风格式满足受 $Co_{\max}$ 控制的正性约束；SLOW 的显式扩散满足
  $Fo_{\max}$ 控制。源码本身不证明整套耦合算法的全局阶数或无条件稳定性。
  \item 刚体墙面规则只是运动学防穿透：整体锐界位置夹在物理墙内，并把指向墙外的对应质心速度
  分量截为零。它不解析接触力、摩擦、恢复或润滑膜，也不改变切向速度和角速度。
  \item 冰水接触 stencil 及阈值 $f_{\rm th}=0.5$ 是当前离散设计参数，不应提升为与网格无关的
  材料常数。
  \item 当前只融化原始材料冰，脱离的水不能重新冻结；空气没有热输运，$90\,^{\circ}\mathrm C$
  水也不蒸发或沸腾。
  \item 剩余固冰必须仍可视为一个连通刚体。若侵蚀后分裂成多个部分，单质心和单角速度模型失效。
  \item sharp 掩膜、连续覆盖率、相场和热水 aperture 是四个相关但不相同的几何描述；用一个阈值
  同时替代四者会破坏质量、能量或动量闭合。
\end{itemize}

\section{运行、输出和阅读结果}

完整当前默认运行：
\begin{lstlisting}[style=iceflow]
python -m iceflow2d
\end{lstlisting}

用于快速验证调度和输出的短运行：
\begin{lstlisting}[style=iceflow]
python -m iceflow2d \
  --resolution-x 100 --resolution-y 200 \
  --end-time-s 0.00005 --output-interval-s 0.00005 \
  --no-plot --save-npz --no-progress
\end{lstlisting}

默认输出目录为
\path{outputs/iceflow2d/coupled_falling_ice_melting_2d/}。始终写出
\path{history.csv} 和 \path{metadata.json}；绘图启用时写出速度、涡量、温度 PNG/GIF 及终态图；
只有显式传入 \code{--save-npz} 才写 \path{fields.npz}。流式图像路径只保留当前帧，峰值内存为
$O(N_xN_y)$；NPZ 路径保留全部采样场，内存为 $O(N_fN_xN_y)$。

阅读输出时应使用半力速度 $\vect u^\star$；世界温度和体积焓是由守恒广延量重建的显示场。
若观察到水--气界面温度颜色快速闪烁，首先要分辨那是“空气固定显示温度与部分充水格温度之间的
阈值切换”，还是 $S^w/V^w$ 真正发生非物理振荡。配对 $V^w/S^w$ 对流和 aperture 同步就是为
避免后一种问题而设置，前一种则可能只是显示掩膜在弥散界面上来回切换。

\section{实现步骤与物理守恒的对应关系}

\begin{table}[H]
\centering
\caption{每类操作修改哪些状态、保护哪个不变量}
\begin{tabularx}{\textwidth}{@{}p{0.26\textwidth}YY@{}}
\toprule
操作 & 主要修改 & 设计上保护或诊断的量 \\
\midrule
$f/h$ 碰撞与流送 & $f_q,h_q,\phi,\vect u,p$ & 相场零阶矩、动量矩、压力恢复一致性 \\
cut-link MEM & 反弹 populations，$\vect J_h,L_h$ & 流体--刚体逐链动量交换；静水牵引一致性 \\
刚体积分/墙面投影 & $\vect V_c,\omega,\vect X_c,\theta$ & 当前锐界整体不穿物理墙；只截断穿墙法向平移速度 \\
equilibrium refill & 新启闭节点 $f,h,\phi,p,\vect u$ & 确定性拓扑更新，不读取同核新节点 \\
联合开口重映射 & $V^w,S^w$ & 全局水面积与显热，局部温度外推及极值限制 \\
成对迎风 & $V^w,S^w$ & 两种广延量使用同一通量，均匀温度保持 \\
界面换热/潜热 & $E^s,m^s,V^w,S^w$ & 冰水成对能量、熔化质量与融水体积 \\
相场水量投影 & $\phi,f,h$ & $\sum\phi=V_w^\star$，纯相不动点 \\
融水三矩 lift & $f,\vect u$ & 两个线动量矩与一个角动量矩 \\
aperture 同步 & $V^w,S^w,T,H$ & 固定总 $V^w,S^w$ 下的合法容量 \\
\bottomrule
\end{tabularx}
\end{table}

\section{结语}

IceFlow2D 的关键并非把多个求解器简单串联，而是在同一步内明确每一种状态由谁负责：LBM 负责
空气--水相场和流体动量；锐 cut-link 负责流体--冰冲量；刚体积分负责侵蚀后的质量、质心和惯量推进，
整体位置与法向速度投影只负责防止冰体穿过物理墙；
材料网格负责固冰质量和显热；世界热网格负责液水面积和显热。相场水量投影、融水三矩 lift
和联合 aperture 重映射分别作用于不同未知量，不能相互替代。严格保持
算法~\ref{alg:step} 的顺序，才可使新位姿、新锐界、新温度、新水量和新动量在 $n+1$ 层重新一致。

\begin{thebibliography}{9}
\bibitem{liang2018}
H. Liang, J. Xu, J. Chen, H. Wang, Z. Chai, and B. Shi,
\emph{Phase-field-based lattice Boltzmann modeling of large-density-ratio two-phase flows},
\emph{Physical Review E}, vol.~97, 033309, 2018.

\bibitem{tao2018}
S. Tao, Q. He, B. Chen, X. Yang, and S. Huang,
\emph{One-point second-order curved boundary condition for lattice Boltzmann simulation of suspended particles},
\emph{Computers \& Mathematics with Applications}, vol.~76, pp.~1593--1607, 2018,
doi: 10.1016/j.camwa.2018.07.013.

\bibitem{werner2021}
L. Werner, C. Rettinger, and U. R{\"u}de,
\emph{Coupling fully resolved light particles with the lattice Boltzmann method on adaptively refined grids},
\emph{International Journal for Numerical Methods in Fluids},
vol.~93, no.~11, pp.~3280--3303, 2021, doi: 10.1002/fld.5034.
\end{thebibliography}

\end{document}
